\documentclass[11pt]{article}
\usepackage[a4paper,margin=1in]{geometry}
\usepackage{fontspec}
\usepackage[boldfont=false]{xeCJK}
\setCJKmainfont{FandolSong-Regular.otf}
\usepackage{amsmath}
\usepackage{amssymb}
\usepackage{array}
\usepackage{booktabs}
\usepackage{graphicx}
\usepackage{adjustbox}
\usepackage{enumitem}
\setlist[itemize]{topsep=2pt,partopsep=0pt,parsep=0pt,itemsep=2pt,leftmargin=1.4em}
\setlist[enumerate]{topsep=2pt,partopsep=0pt,parsep=0pt,itemsep=2pt,leftmargin=1.6em}
\usepackage{parskip}
\usepackage{fvextra}
\fvset{breaklines=true,breakanywhere=true,fontsize=\small}
\setlength{\parindent}{0pt}

\begin{document}

\begin{center}
  {\LARGE\bfseries Biological molecules}\\[0.35em]
  {\large Igcse Biology}
\end{center}
\vspace{0.6em}

\section*{The chemical elements in food}

Living things are built from a few kinds of \textbf{molecule} 分子. The most important are carbohydrates, fats and proteins. Each is made from a small number of chemical \textbf{elements} 元素.

\begin{itemize}
\item \textbf{Carbohydrates} 碳水化合物 and \textbf{fats} 脂肪 contain three elements: \textbf{carbon} 碳 (C), \textbf{hydrogen} 氢 (H) and oxygen (O).

\item \textbf{Proteins} 蛋白质 contain carbon, hydrogen and oxygen too, plus \textbf{nitrogen} 氮 (N). (Some proteins also contain sulfur, S.)

\end{itemize}

\section*{Building large molecules from small ones}

Large molecules are made by joining many small molecules together, like beads on a string.

\begin{center}
\adjustbox{max width=\linewidth}{%
\begin{tabular}{ll}
\toprule
\textbf{Large molecule} & \textbf{Built from} \\
\midrule
\textbf{starch} 淀粉, \textbf{glycogen} 糖原 and \textbf{cellulose} 纤维素 & many \textbf{glucose} 葡萄糖 units \\
proteins & \textbf{amino acids} 氨基酸 \\
fats and oils & \textbf{fatty acids} 脂肪酸 and \textbf{glycerol} 甘油 \\
\bottomrule
\end{tabular}}
\end{center}

Starch and glycogen are energy stores; cellulose makes plant cell walls. Proteins are built from about 20 different kinds of amino acid joined in a chain. One fat molecule is made from three fatty acids joined to one glycerol.

\section*{Food tests}

You can test a piece of food to find out which substances it contains. Add the test chemical and watch for a colour change.

\begin{center}
\adjustbox{max width=\linewidth}{%
\begin{tabular}{lll}
\toprule
\textbf{Food substance} & \textbf{Test} & \textbf{Positive result} \\
\midrule
starch & add \textbf{iodine solution} 碘液 & orange-brown → blue-black \\
\textbf{reducing sugars} 还原糖 (such as glucose) & add \textbf{Benedict's solution} 本尼迪特试剂 and heat & blue → brick-red / orange \\
protein & add \textbf{biuret} 双缩脲 solution & blue → purple \\
fats and oils & \textbf{ethanol} 乙醇 \textbf{emulsion} 乳浊液 test & a cloudy white layer forms \\
\textbf{vitamin} 维生素 C & add to blue DCPIP & blue → colourless \\
\bottomrule
\end{tabular}}
\end{center}

Two things to remember: only the Benedict's test needs \textbf{heat}; and DCPIP loses its colour (goes colourless), while the other tests gain a new colour.

\section*{The structure of DNA (Supplement)}

\textbf{DNA} is the \textbf{genetic material} 遗传物质 — it carries the instructions for building and running an organism. Its structure has four key points:

\begin{itemize}
\item two \textbf{strands} 链 are coiled together to form a \textbf{double helix} 双螺旋 (like a twisted ladder).

\item each strand carries chemicals called \textbf{bases} 碱基.

\item \textbf{bonds} 化学键 between pairs of bases hold the two strands together.

\item the bases always pair in the same way: \textbf{A with T}, and \textbf{C with G}.

\end{itemize}

\section*{Exam tips}

\begin{itemize}
\item Carbohydrates and fats are made of C, H and O. Proteins also contain N (nitrogen).

\item Learn each food test and its colour change. For reducing sugars you must \textbf{heat} the Benedict's solution.

\item DCPIP goes from blue to \textbf{colourless}; the other four tests each give a \textbf{new} colour.

\item In DNA, the bases pair only A–T and C–G. So if you know the bases on one strand, you can work out the other.

\end{itemize}


\end{document}
